\section{Related Work}
\label{sec:related_work}

\subsection{Machine Unlearning for LLMs}

The concept of machine unlearning was first formalized by \citet{cao2015towards}, who proposed efficient data removal from learned models without full retraining. \citet{bourtoule2021machine} later introduced SISA training, which partitions data during training to make unlearning tractable. However, these approaches target classification models and do not apply to large autoregressive language models.

For LLMs, gradient ascent (GA) is the most direct approach: maximize the loss on the forget set to push the model away from memorized content \citep{jang2023knowledge}. GA is unbounded, the loss grows without limit, and it rapidly destroys model coherence on all inputs. \citet{yao2023large} proposed GradDiff, combining gradient ascent on forget data with cross-entropy on retain data. But retain performance still degrades because the forget loss drives weights in destructive directions.

A preference-optimization line of work began with NPO \citep{zhang2024negative}, which adapts DPO \citep{rafailov2023direct} by treating forget samples as negative preferences. NPO collapses exponentially slower than GA, but it requires a frozen reference model in memory and its objective saturates instead of converging to a well-defined target. \citet{fan2024simplicity} address the reference-model cost with SimNPO, dropping the reference model entirely while remaining competitive on TOFU and MUSE. \citet{bu2024unlearning} take a multi-task view with NGDiff, normalizing the gradient difference between forget and retain objectives so that one does not dominate the other.

At the representation level, \citet{li2024wmdp} proposed RMU as part of the WMDP benchmark. RMU steers intermediate representations toward random targets for forget data while anchoring retain representations. It works well in their benchmark but lacks the optimization-theoretic grounding of loss-based methods.

What all these approaches have in common is the absence of a \emph{bounded, self-stabilizing} forget objective with a per-token stopping criterion. GA is unbounded and diverges. NPO's adaptive weighting damps asymptotically but never reaches exactly zero---it still applies gradient to well-forgotten sequences. Na\"ive entropy maximization \citep{yuan2024closer} pushes \emph{all} tokens toward uniform indefinitely, with no concept of ``sufficiently forgotten.'' The inverted hinge loss of LoKU \citep{cha2024loku} is bounded but operates on pairwise token margins, not distributional entropy, and provides no tunable threshold.

Our clamped entropy loss is qualitatively different: it defines a vocabulary-invariant sufficiency threshold ($\tau \cdot \log V$) and produces \emph{exactly zero} gradient once a token crosses it. This hard deadzone means forgotten tokens permanently drop out of optimization---the loss is bounded by construction, the gradient norm shrinks monotonically, and no reference model is needed.

\subsection{Unlearning Benchmarks and Evaluation}

Evaluating unlearning is arguably as hard as the unlearning itself. \citet{maini2024tofu} introduced TOFU, a benchmark built around 200 synthetic author biographies with associated QA pairs. Since the data is fictitious, a gold-standard retrained model naturally exists. They evaluate using truth ratios, ROUGE scores, and membership inference attacks, and find that no tested baseline achieves truly effective unlearning.

\citet{shi2024muse} proposed MUSE, which defines six desirable properties for an unlearned model: no verbatim memorization, no knowledge memorization, no privacy leakage, utility preservation, scalability, and sustainability. They evaluate on two realistic corpora (news articles and Harry Potter books) with Llama-2-7b as the base model. Their results show that while some methods reduce memorization, most fail on privacy leakage or severely degrade utility. The WMDP benchmark \citep{li2024wmdp} takes a safety-oriented perspective, measuring hazardous knowledge in biosecurity and cybersecurity domains.

We evaluate on both TOFU and MUSE because they provide complementary views: TOFU tests unlearning of factual associations in a controlled setting, while MUSE tests removal of naturalistic memorized content at larger scale.

\subsection{Parameter-Efficient Unlearning}

LoRA \citep{hu2022lora} has become the standard method for parameter-efficient adaptation of LLMs. It factorizes weight updates into low-rank matrices, reducing trainable parameters and memory by orders of magnitude. Several recent works apply LoRA to unlearning. \citet{cha2024loku} proposed LoKU, which addresses the poor forget--retain trade-off when combining GA with LoRA by introducing an inverted hinge loss and Fisher-informed LoRA initialization. \citet{jia2024wagle} proposed WAGLE, which assigns importance scores to model weights and applies strategic masking for modular unlearning.

In \method, we use LoRA not just for memory efficiency but as a \emph{structural constraint}: the low-rank bottleneck limits the expressivity of the unlearning update, preventing the catastrophic weight changes that happen with full-parameter methods. This complements our bilevel formulation. The inner loop protects retain performance within a small parameter subspace, while the outer loop drives forgetting through the same bottleneck.

\subsection{Bilevel Optimization}

Bilevel optimization has a long history in operations research \citep{bertsekas1999nonlinear} and is increasingly used in machine learning. MAML \citep{finn2017model} optimizes for fast adaptation in the outer loop while performing task-specific fine-tuning in the inner loop. \citet{franceschi2018bilevel} provided a unified bilevel framework for hyperparameter optimization and meta-learning. Computing exact hypergradients requires second-order information, so various approximations exist: implicit differentiation \citep{lorraine2020optimizing}, Neumann series for inverse Hessian approximation \citep{ji2021bilevel}, and first-order approaches that avoid Hessian computation entirely.

The augmented Lagrangian method (ALM) \citep{bertsekas1999nonlinear, nocedal2006numerical} maintains dual variables alongside the primal objective for constrained optimization. Unlike penalty methods that require driving the penalty to infinity, ALM can achieve exact constraint satisfaction with finite parameters through adaptive dual updates. ALM is well-suited for our setting because retain performance must be strictly enforced as a constraint, not loosely balanced as a loss term.

The most closely related works along the constrained optimization axis are BLUR \citep{reisizadeh2025blur}, PDU \citep{entesari2025constrained}, and \citet{cheng2026unlearning}. BLUR uses the retain loss as the lower-level objective and the forget loss as the upper-level one, optimizing through implicit differentiation. PDU uses a standard Lagrangian (linear penalty only, no quadratic term) with a logit-margin forget loss and symmetric dual updates. \citet{cheng2026unlearning} use an augmented Lagrangian for unlearning, but with an \emph{equality} constraint ($L_{\text{ret}} = L_{\text{ret}}(\thetav_0)$, penalty always active), a flat single-level optimizer (not bilevel), and evaluation only on vision models (ResNet-18, ViT-B) and classification tasks.

Our approach differs from all three in structural ways: (a) we enforce an \emph{inequality} constraint $L_{\text{ret}} \leq \varepsilon$ where the quadratic penalty activates only on violation---the model is never penalized for being better than baseline, unlike equality-constrained formulations; (b) we use an \emph{asymmetric} dual update (fast $\lambda$ increase on violation, 10$\times$ slower decay when satisfied), producing one-shot ratchet behavior rather than the recurring oscillations of symmetric updates; (c) the retain objective is nested in the \emph{inner loop} of a bilevel problem, creating a ``repair first, forget carefully'' dynamic that inverts SCRUB's inner-loop forgetting; and (d) our forget objective is clamped entropy, a bounded loss with a hard per-token deadzone---in contrast to PDU's unbounded logit margin and the unbounded gradient ascent used in BLUR. \citet{jia2024soul} also explored second-order optimization for unlearning through Newton-step approximations via influence functions \citep{koh2017understanding}.

Our earlier attempt at bilevel unlearning used binary sparsity masks to select which parameters to update. LoRA's continuous low-rank structure turned out to be both more stable and more effective, which motivated the current \method design.
